设 $n\ge 2$、$1\le p\le\infty$，且 $f\in L^p(\mathbb R^n)$ 任意。对于 $r>0$，令 $B_r^\alpha$ 为作用在 $\mathbb R^n$ 上、符号为 $(1-|\xi|^2/r^2)_+^\alpha$ 的 Fourier 乘子，并令
\[
\delta(p)=\max\{0,n|1/p-1/2|-1/2\}.
\]
若 $\alpha>\delta(p)$，则
\[
\sup_{r>0}\|B_r^\alpha f\|_{L^p(\mathbb R^n)}
\le C_{n,p,\alpha}\|f\|_{L^p(\mathbb R^n)}.
\]
在对 $\alpha$ 的相同条件下，平均值收敛到 $f$，收敛发生于 $L^p$ 中、在 $r\to\infty$ 时成立，其中要求 $p<\infty$。若 $2\le p<\infty$ 且 $\alpha>\max\{0,(n-1)/2-n/p\}$，则还有
\[
\left\|\sup_{r>0}|B_r^\alpha f|\right\|_{L^p(\mathbb R^n)}
\le C_{n,p,\alpha}\|f\|_{L^p(\mathbb R^n)},
\]
且这些平均值几乎处处收敛到 $f$，此时 $r\to\infty$。上确界针对一个几乎处处关于 $r>0$ 连续的代表元取值。
